% GENERATED FILE. Edit canonical lessons, then run npm run generate:books.
% canonical-source-hash: fnv1a64:8a942f358928e554
% canonical-lessons: SA-C38-fire, SA-C38-ash, SA-C38-stone, SA-C38-sand, SA-C38-soil

\chapter{Fire and Stone}
\label{ch:sa-fire-and-stone}
\hlchaptermodality{38}

\begin{chapteropening}
\textbf{By the end of this chapter:} \emph{I can name fire, ash, a stone, sand and soil in Sanskrit.}

Complete SA-C38-soil and demonstrate the chapter promise.
\end{chapteropening}

\section[agniḥ]{\sk{अग्निः} (agniḥ) — fire}

\label{lesson:SA-C38-fire}



\begin{quote}
Before the new word: what did \emph{āśīrvādaḥ} mean?
\end{quote}

\subsection*{You'll want to know: \sk{अग्निः}}
\textbf{\sk{अग्निः}} (\emph{agniḥ}) — \textquotedblleft{}fire\textquotedblright{}.

\textbf{\sk{अग्निः}} is fire itself — the flame on a hearth, and the fire an offering is poured into.

You have had the doing-word for a while: \sk{दहति}, \textquotedblleft{}burns\textquotedblright{}. \textbf{\sk{अग्निः}} is the thing that does it.

Latin \emph{ignis} is the very same word, and English carries it in \emph{ignite} and \emph{igneous}; Russian \emph{ogon'} is another of its brothers. It went west almost unchanged, and in Sanskrit it stands as the first word of the oldest hymn there is.

\begin{grammarlens}[title={What you've built}]
Fire, named as a thing rather than as an action.
\end{grammarlens}

\subsection*{Your turn}
Say these aloud:

\begin{itemize}
\raggedright
  \item \emph{agniḥ}
  \item it once more, slowly
  \item \emph{āśīrvādaḥ}, then \emph{agniḥ}
\end{itemize}

\begin{tcolorbox}[breakable,colback=teal!4,colframe=teal!35!black,title={Before you move on}]
What does \sk{अग्निः} mean? (\textquotedblleft{}fire\textquotedblright{}.) Say it once more.
\end{tcolorbox}
\section[bhasma]{\sk{भस्म} (bhasma) — ash}

\label{lesson:SA-C38-ash}



\begin{quote}
Before the new word: what did \emph{agniḥ} mean?
\end{quote}

\subsection*{You'll want to know: \sk{भस्म}}
\textbf{\sk{भस्म}} (\emph{bhasma}) — \textquotedblleft{}ash\textquotedblright{}.

\textbf{\sk{भस्म}} is what \sk{अग्निः} leaves behind — ash, and the grey line of it drawn across a forehead.

It stands on \emph{bhas}, \textquotedblleft{}to chew, to grind small\textquotedblright{}, so ash is the ground-down thing: what is left when fire has finished eating.

Hindi and Marathi keep \emph{bhasm} with the meaning untouched.

\begin{grammarlens}[title={What you've built}]
Fire, and what fire leaves.
\end{grammarlens}

\subsection*{Your turn}
Say these aloud:

\begin{itemize}
\raggedright
  \item \emph{bhasma}
  \item it once more, slowly
  \item \emph{agniḥ}, then \emph{bhasma}
\end{itemize}

\begin{tcolorbox}[breakable,colback=teal!4,colframe=teal!35!black,title={Before you move on}]
What does \sk{भस्म} mean? (\textquotedblleft{}ash\textquotedblright{}.) Say it once more.
\end{tcolorbox}
\section[śilā]{\sk{शिला} (śilā) — a stone, a slab of rock}

\label{lesson:SA-C38-stone}



\begin{quote}
Before the new word: what did \emph{bhasma} mean?
\end{quote}

\subsection*{You'll want to know: \sk{शिला}}
\textbf{\sk{शिला}} (\emph{śilā}) — \textquotedblleft{}a stone, a slab of rock\textquotedblright{}.

\textbf{\sk{शिला}} is a stone large enough to sit on: a slab, a flat rock, the piece a step is cut from.

Beside \sk{पर्वतः}, which is the whole raised land, \textbf{\sk{शिला}} is one piece of it you can put a hand on.

Its history outside India is unsettled — nobody has agreed a cousin for it, the way nobody has for \sk{केशः} or \sk{माला}. Feminine, ending in \textbf{-\sk{आ}}.

\begin{grammarlens}[title={What you've built}]
One piece of the mountain, small enough to touch.
\end{grammarlens}

\subsection*{Your turn}
Say these aloud:

\begin{itemize}
\raggedright
  \item \emph{śilā}
  \item it once more, slowly
  \item \emph{bhasma}, then \emph{śilā}
\end{itemize}

\begin{tcolorbox}[breakable,colback=teal!4,colframe=teal!35!black,title={Before you move on}]
What does \sk{शिला} mean? (\textquotedblleft{}a stone, a slab of rock\textquotedblright{}.) Say it once more.
\end{tcolorbox}
\section[sikatā]{\sk{सिकता} (sikatā) — sand}

\label{lesson:SA-C38-sand}



\begin{quote}
Before the new word: what did \emph{śilā} mean?
\end{quote}

\subsection*{You'll want to know: \sk{सिकता}}
\textbf{\sk{सिकता}} (\emph{sikatā}) — \textquotedblleft{}sand\textquotedblright{}.

\textbf{\sk{सिकता}} is sand — the grit of a bank, and the bank itself where it piles up.

It has long been tied to \emph{sic}, \textquotedblleft{}to pour, to sprinkle\textquotedblright{}: sand is the stuff that pours.

Set it beside \sk{नदी} and you have the two halves of a bank — the water, and what the water leaves at its edge.

\begin{grammarlens}[title={What you've built}]
What a river puts down at its edge.
\end{grammarlens}

\subsection*{Your turn}
Say these aloud:

\begin{itemize}
\raggedright
  \item \emph{sikatā}
  \item it once more, slowly
  \item \emph{śilā}, then \emph{sikatā}
\end{itemize}

\begin{tcolorbox}[breakable,colback=teal!4,colframe=teal!35!black,title={Before you move on}]
What does \sk{सिकता} mean? (\textquotedblleft{}sand\textquotedblright{}.) Say it once more.
\end{tcolorbox}
\section[mṛttikā]{\sk{मृत्तिका} (mṛttikā) — soil}

\label{lesson:SA-C38-soil}



\begin{quote}
Before the last word of the chapter: what did \emph{śilā} mean, and what did \emph{sikatā} mean?
\end{quote}

\subsection*{You'll want to know: \sk{मृत्तिका}}
\textbf{\sk{मृत्तिका}} (\emph{mṛttikā}) — \textquotedblleft{}soil\textquotedblright{}.

\textbf{\sk{मृत्तिका}} is soil — damp, workable, the stuff a potter digs and a wall is plastered with.

It stands on \emph{mṛd}, and Latin \emph{mordēre}, \textquotedblleft{}to bite\textquotedblright{}, is its cousin: the sense underneath both is pressing something soft until it gives. English took \emph{mordant} from that Latin side.

Eastward it wore down into Hindi \emph{miṭṭī}, still the ordinary word for the ground you dig.

\begin{grammarlens}[title={What you've built}]
Five things underfoot and in the hearth: fire, ash, stone, sand, soil.
\end{grammarlens}

\subsection*{Your turn}
Say these aloud:

\begin{itemize}
\raggedright
  \item \emph{mṛttikā}
  \item it once more, slowly
  \item all five in order, then \emph{agniḥ} and \emph{mṛttikā} together
\end{itemize}

\begin{tcolorbox}[breakable,colback=teal!4,colframe=teal!35!black,title={Before you move on}]
What does \sk{मृत्तिका} mean? (\textquotedblleft{}soil\textquotedblright{}.) Say it once more.
\end{tcolorbox}
